% ENEM 2016-2025 — 17 questões MCQ — Porcentagem
% Fonte: lista "Porcentagem I" do site projetoagathaedu.com.br
% Omitidas: Q3,Q5,Q8,Q13,Q15,Q16,Q17,Q21,Q22,Q25,Q27,Q31 (imagens)
%           Q12,Q14 (duplicatas já no banco)
% Nota Q9: gabarito do site (D) está errado; correto é B (1% menor) — verificado por cálculo

---
tipo: Múltipla Escolha
dificuldade: Média
nivel: médio
disciplina: Matemática
assunto: Porcentagem
gabarito: A
tags: [ENEM, porcentagem, desconto, papel, proporcionalidade]
fonte: concurso
concurso: ENEM
banca: INEP
ano: 2025
numero: 1
---
\question
A cada bimestre, a diretora de uma escola compra uma quantidade de folhas de papel ofício proporcional ao número de alunos matriculados. No bimestre passado, ela comprou 6.000 folhas para serem utilizadas pelos 1.200 alunos matriculados. Neste bimestre, alguns alunos cancelaram suas matrículas e a escola tem, agora, 1.150 alunos. A diretora só pode gastar R\$ 220,00 nessa compra, e sabe que o fornecedor da escola vende as folhas de papel ofício em embalagens de 100 unidades a R\$ 4,00 a embalagem. Assim, será preciso convencer o fornecedor a dar um desconto à escola, de modo que seja possível comprar a quantidade total de papel ofício necessária para o bimestre.

O desconto necessário no preço final da compra, em porcentagem, pertence ao intervalo
\begin{choices}
\CorrectChoice (5,0 ; 5,5).
\choice (8,0 ; 8,5).
\choice (11,5 ; 12,5).
\choice (19,5 ; 20,5).
\choice (3,5 ; 4,0).
\end{choices}

---
tipo: Múltipla Escolha
dificuldade: Média
nivel: médio
disciplina: Matemática
assunto: Porcentagem
gabarito: B
tags: [ENEM, porcentagem, creche, lista de espera, redução]
fonte: concurso
concurso: ENEM
banca: INEP
ano: 2025
numero: 2
---
\question
Em janeiro do ano passado, a direção de uma fábrica abriu uma creche para os filhos de seus funcionários, com 10 salas, cada uma com capacidade para atender 10 crianças a cada ano. As vagas são sorteadas entre os filhos dos funcionários inscritos, enquanto os não contemplados pelo sorteio ficam numa lista de espera. No ano seguinte, a lista de espera teve 400 nomes e, neste ano, esse número cresceu 10\%.

A direção da fábrica realizou uma pesquisa e constatou que a lista de espera para o próximo ano terá a mesma quantidade de nomes da lista de espera deste ano. Decidiu, então, construir, ao longo desse ano, novas salas para a creche, também com capacidade de atendimento para 10 crianças cada, de modo que o número de nomes na lista de espera no próximo ano seja 25\% menor que o deste ano.

O número mínimo de salas que deverão ser construídas é
\begin{choices}
\choice 10.
\CorrectChoice 11.
\choice 13.
\choice 30.
\choice 33.
\end{choices}

---
tipo: Múltipla Escolha
dificuldade: Média
nivel: médio
disciplina: Matemática
assunto: Porcentagem
gabarito: D
tags: [ENEM, porcentagem, erro de medição, bagagem, balança]
fonte: concurso
concurso: ENEM
banca: INEP
ano: 2022
cargo: PPL
numero: 4
---
\question
Os aeroportos apresentam regras rígidas para despacho de bagagem. No caso de embarque nacional, algumas companhias aéreas ainda não cobravam, até 2017, por unidade de bagagem despachada, limitada a 23 kg por passageiro.

Uma pessoa irá viajar com uma única mala. Como não quer pagar por ``excesso de peso'' e dispõe, em casa, de uma balança de pêndulo que apresenta um erro máximo de 8\% a mais em relação à massa real do objeto que nela for verificada, conferirá qual a massa de sua mala antes de ir para o aeroporto.

O valor máximo, em quilograma, indicado em sua balança deverá ser
\begin{choices}
\choice 21,16.
\choice 22,08.
\choice 23,92.
\CorrectChoice 24,84.
\choice 25,00.
\end{choices}

---
tipo: Múltipla Escolha
dificuldade: Fácil
nivel: médio
disciplina: Matemática
assunto: Porcentagem
gabarito: D
tags: [ENEM, porcentagem, fração, língua estrangeira]
fonte: concurso
concurso: ENEM
banca: INEP
ano: 2022
cargo: PPL
numero: 6
---
\question
Uma escola realizou uma pesquisa entre todos os seus estudantes e constatou que três em cada dez deles estão matriculados em algum curso extracurricular de língua estrangeira.

Em relação ao número total de estudantes dessa escola, qual porcentagem representa o número de alunos matriculados em algum curso extracurricular de língua estrangeira?
\begin{choices}
\choice 0,3\%
\choice 0,33\%
\choice 3\%
\CorrectChoice 30\%
\choice 33\%
\end{choices}

---
tipo: Múltipla Escolha
dificuldade: Média
nivel: médio
disciplina: Matemática
assunto: Porcentagem
gabarito: A
tags: [ENEM, porcentagem, desconto, cartão de crédito, promoção]
fonte: concurso
concurso: ENEM
banca: INEP
ano: 2025
numero: 7
---
\question
Em uma loja, o preço promocional de uma geladeira é de R\$ 1.000,00 para pagamento somente em dinheiro. Seu preço normal, fora da promoção, é 10\% maior. Para pagamento feito com o cartão de crédito da loja, é dado um desconto de 2\% sobre o preço normal.

Uma cliente decidiu comprar essa geladeira, optando pelo pagamento com o cartão de crédito da loja. Ela calculou que o valor a ser pago seria o preço promocional acrescido de 8\%. Ao ser informada pela loja do valor a pagar, segundo sua opção, percebeu uma diferença entre seu cálculo e o valor que lhe foi apresentado.

O valor apresentado pela loja, comparado ao valor calculado pela cliente, foi
\begin{choices}
\CorrectChoice R\$ 2,00 menor.
\choice R\$ 100,00 menor.
\choice R\$ 200,00 menor.
\choice R\$ 42,00 maior.
\choice R\$ 80,00 maior.
\end{choices}

---
tipo: Múltipla Escolha
dificuldade: Média
nivel: médio
disciplina: Matemática
assunto: Porcentagem
gabarito: B
tags: [ENEM, porcentagem, combustível, aeronave, variação percentual]
fonte: concurso
concurso: ENEM
banca: INEP
ano: 2025
numero: 9
---
\question
Para realizar um voo entre duas cidades que distam 2.000 km uma da outra, uma companhia aérea utilizava um modelo de aeronave A, capaz de transportar até 200 passageiros. Quando uma dessas aeronaves está lotada de passageiros, o consumo de combustível é de 0,02 litro por quilômetro e por passageiro. Essa companhia resolveu trocar o modelo de aeronave A pelo modelo de aeronave B, que é capaz de transportar 10\% de passageiros a mais do que o modelo A, mas consumindo 10\% menos combustível por quilômetro e por passageiro.

A quantidade de combustível consumida pelo modelo de aeronave B, em relação à do modelo de aeronave A, em um voo lotado entre as duas cidades, é
\begin{choices}
\choice 10\% menor.
\CorrectChoice 1\% menor.
\choice igual.
\choice 1\% maior.
\choice 11\% maior.
\end{choices}

---
tipo: Múltipla Escolha
dificuldade: Média
nivel: médio
disciplina: Matemática
assunto: Porcentagem
gabarito: B
tags: [ENEM, porcentagem, desconto, reajuste, celular]
fonte: concurso
concurso: ENEM
banca: INEP
ano: 2022
cargo: PPL
numero: 10
---
\question
A associação de comerciantes varejistas de uma cidade, a fim de incrementar as vendas para o Natal, decidiu promover um fim de semana de descontos e promoções, no qual produtos e serviços estariam com valores reduzidos. Antes do período promocional, um celular custava R\$ 300,00 e teve seu preço reajustado, passando a custar R\$ 315,00. Durante o fim de semana de descontos e promoções, o preço desse celular recebeu um desconto de 20\%.

O desconto dado no preço do celular, em porcentagem, com base no valor dele anteriormente ao aumento sofrido antes da promoção, foi de
\begin{choices}
\choice 15,24\%
\CorrectChoice 16,00\%
\choice 19,04\%
\choice 21,00\%
\choice 25,00\%
\end{choices}

---
tipo: Múltipla Escolha
dificuldade: Difícil
nivel: médio
disciplina: Matemática
assunto: Porcentagem
gabarito: C
tags: [ENEM, porcentagem, marketing digital, redes sociais, proporcionalidade]
fonte: concurso
concurso: ENEM
banca: INEP
ano: 2025
numero: 11
---
\question
Uma equipe de marketing digital foi contratada para aumentar as vendas de um produto ofertado em um site de comércio eletrônico. Para isso, elaborou um anúncio que, quando o cliente clica sobre ele, é direcionado para a página de vendas do produto. Esse anúncio foi divulgado em duas redes sociais, A e B, e foram obtidos os seguintes resultados:

\begin{itemize}
\item rede social A: o anúncio foi visualizado por 3.000 pessoas; 10\% delas clicaram sobre o anúncio e foram redirecionadas para o site; 3\% das que clicaram sobre o anúncio compraram o produto. O investimento feito para a publicação do anúncio nessa rede foi de R\$ 100,00;
\item rede social B: o anúncio foi visualizado por 1.000 pessoas; 30\% delas clicaram sobre o anúncio e foram redirecionadas para o site; 2\% das que clicaram sobre o anúncio compraram o produto. O investimento feito para a publicação do anúncio nessa rede foi de R\$ 200,00.
\end{itemize}

Por experiência, o pessoal da equipe de marketing considera que a quantidade de novas pessoas que verão o anúncio é diretamente proporcional ao investimento realizado, e que a quantidade de pessoas que comprarão o produto também se manterá proporcional à quantidade de pessoas que clicarão sobre o anúncio.

O responsável pelo produto decidiu, então, investir mais R\$ 300,00 em cada uma das duas redes sociais para a divulgação desse anúncio e obteve, de fato, o aumento proporcional esperado na quantidade de clientes que compraram esse produto. Para classificar o aumento obtido na quantidade (Q) de compradores desse produto, em consequência dessa segunda divulgação, em relação aos resultados observados na primeira divulgação, o responsável pelo produto adotou o seguinte critério:

\begin{itemize}
\item $Q \leq 60\%$: não satisfatório;
\item $60\% < Q \leq 100\%$: regular;
\item $100\% < Q \leq 150\%$: bom;
\item $150\% < Q \leq 190\%$: muito bom;
\item $190\% < Q \leq 200\%$: excelente.
\end{itemize}

O aumento na quantidade de compradores, em consequência dessa segunda divulgação, em relação ao que foi registrado com a primeira divulgação, foi classificado como
\begin{choices}
\choice não satisfatório.
\choice regular.
\CorrectChoice bom.
\choice muito bom.
\choice excelente.
\end{choices}

---
tipo: Múltipla Escolha
dificuldade: Média
nivel: médio
disciplina: Matemática
assunto: Porcentagem
gabarito: A
tags: [ENEM, porcentagem, reajuste, editora, papel]
fonte: concurso
concurso: ENEM
banca: INEP
ano: 2020
cargo: Digital
numero: 18
---
\question
Uma editora pretende fazer uma reimpressão de um de seus livros. A direção da editora sabe que o gasto com papel representa 60\% do custo de reimpressão, e que as despesas com a gráfica representam os 40\% restantes. Dentro da programação da editora, no momento em que ela realizar a reimpressão, o preço do papel e os custos com a gráfica terão sofrido reajustes de 25,9\% e 32,5\%, respectivamente. O custo para a reimpressão de cada livro, nos preços atuais, é de R\$ 100,00.

Qual será o custo, em real, para a reimpressão de cada livro com os reajustes estimados de custo de papel e despesas com a gráfica?
\begin{choices}
\CorrectChoice 128,54
\choice 129,20
\choice 129,86
\choice 158,40
\choice 166,82
\end{choices}

---
tipo: Múltipla Escolha
dificuldade: Média
nivel: médio
disciplina: Matemática
assunto: Porcentagem
gabarito: C
tags: [ENEM, porcentagem, área, lote, lazer]
fonte: concurso
concurso: ENEM
banca: INEP
ano: 2020
cargo: Digital
numero: 19
---
\question
Uma pessoa possuía um lote com área de 300 m². Nele construiu sua casa, utilizando 70\% do lote para construção da residência e o restante para área de lazer. Posteriormente, adquiriu um novo lote ao lado do de sua casa e, com isso, passou a dispor de um terreno formado pelos dois lotes, cuja área mede 420 m². Decidiu então ampliar a casa, de tal forma que ela ocupasse no mínimo 60\% da área do terreno, sendo o restante destinado à área de lazer.

O acréscimo máximo que a região a ser destinada à área de lazer no terreno poderá ter, em relação à área que fora utilizada para lazer no lote original, em metro quadrado, é
\begin{choices}
\choice 12
\choice 48
\CorrectChoice 78
\choice 138
\choice 168
\end{choices}

---
tipo: Múltipla Escolha
dificuldade: Difícil
nivel: médio
disciplina: Matemática
assunto: Porcentagem
gabarito: E
tags: [ENEM, porcentagem, IDEB, variação percentual, diretamente proporcional, inversamente proporcional]
fonte: concurso
concurso: ENEM
banca: INEP
ano: 2020
cargo: Digital
numero: 20
---
\question
O Índice de Desenvolvimento da Educação Básica (Ideb), criado para medir a qualidade do aprendizado do ensino básico no Brasil, é calculado a cada dois anos. No seu cálculo são combinados dois indicadores: o aprendizado e o fluxo escolar.

O Ideb de uma escola numa dada série escolar pode ser calculado pela expressão $\text{Ideb} = N \times P$, em que $N$ é a média da proficiência em língua portuguesa e matemática, variando de 0 a 10, e $P = \dfrac{1}{T}$, em que $T$ é o tempo médio de permanência dos alunos na série.

Uma escola apresentou no 9º ano do ensino fundamental, em 2017, um Ideb diferente daquele que havia apresentado nessa mesma série em 2015, pois o tempo médio de permanência dos alunos no 9º ano diminuiu 2\%, enquanto a média de proficiência em língua portuguesa e matemática, nessa série, aumentou em 2\%.

Dessa forma, o Ideb do 9º ano do ensino fundamental dessa escola em 2017, em relação ao calculado em 2015,
\begin{choices}
\choice permaneceu inalterado, pois o aumento e a diminuição de 2\% nos dois parâmetros anulam-se.
\choice aumentou em 4\%, pois o aumento de 2\% na média da proficiência soma-se à diminuição de 2\% no tempo médio de permanência dos alunos na série.
\choice diminuiu em 4,04\%, pois tanto o decrescimento do tempo médio de permanência dos alunos na série em 2\% quanto o crescimento da média da proficiência em 2\% implicam dois decréscimos consecutivos de 2\% no valor do Ideb.
\choice aumentou em 4,04\%, pois tanto o decrescimento do tempo médio de permanência dos alunos na série em 2\% quanto o crescimento da média da proficiência em 2\% implicam dois acréscimos consecutivos de 2\% no valor do Ideb.
\CorrectChoice aumentou em 4,08\%, pois houve um acréscimo de 2\% num parâmetro que é diretamente proporcional e um decréscimo de 2\% num parâmetro que é inversamente proporcional ao Ideb.
\end{choices}

---
tipo: Múltipla Escolha
dificuldade: Média
nivel: médio
disciplina: Matemática
assunto: Porcentagem
gabarito: E
tags: [ENEM, porcentagem, telefone, redução, conta]
fonte: concurso
concurso: ENEM
banca: INEP
ano: 2019
cargo: PPL
numero: 23
---
\question
A conta de telefone de uma loja foi, nesse mês, de R\$ 200,00. O valor da assinatura mensal, já incluso na conta, é de R\$ 40,00, o qual dá direito a realizar uma quantidade ilimitada de ligações locais para telefones fixos. As ligações para celulares são tarifadas separadamente. Nessa loja, são feitas somente ligações locais, tanto para telefones fixos quanto para celulares. Para reduzir os custos, o gerente planeja, para o próximo mês, uma conta de telefone com valor de R\$ 80,00.

Para que esse planejamento se cumpra, a redução percentual com gastos em ligações para celulares nessa loja deverá ser de
\begin{choices}
\choice 25\%
\choice 40\%
\choice 50\%
\choice 60\%
\CorrectChoice 75\%
\end{choices}

---
tipo: Múltipla Escolha
dificuldade: Média
nivel: médio
disciplina: Matemática
assunto: Porcentagem
gabarito: E
tags: [ENEM, porcentagem, desconto, frete, menor preço]
fonte: concurso
concurso: ENEM
banca: INEP
ano: 2019
cargo: PPL
numero: 24
---
\question
Deseja-se comprar determinado produto e, após uma pesquisa de preços, o produto foi encontrado em 5 lojas diferentes, a preços variados.

\begin{itemize}
\item Loja 1: 20\% de desconto, que equivale a R\$ 720,00, mais R\$ 70,00 de frete;
\item Loja 2: 20\% de desconto, que equivale a R\$ 740,00, mais R\$ 50,00 de frete;
\item Loja 3: 20\% de desconto, que equivale a R\$ 760,00, mais R\$ 80,00 de frete;
\item Loja 4: 15\% de desconto, que equivale a R\$ 710,00, mais R\$ 10,00 de frete;
\item Loja 5: 15\% de desconto, que equivale a R\$ 690,00, sem custo de frete.
\end{itemize}

O produto foi comprado na loja que apresentou o menor preço total.

O produto foi adquirido na loja
\begin{choices}
\choice 1.
\choice 2.
\choice 3.
\choice 4.
\CorrectChoice 5.
\end{choices}

---
tipo: Múltipla Escolha
dificuldade: Média
nivel: médio
disciplina: Matemática
assunto: Porcentagem
gabarito: C
tags: [ENEM, menor custo, mirante, elevador, teleférico]
fonte: concurso
concurso: ENEM
banca: INEP
ano: 2025
numero: 26
---
\question
Em um parque há dois mirantes de alturas distintas que são acessados por elevador panorâmico. O topo do mirante 1 é acessado pelo elevador 1, enquanto que o topo do mirante 2 é acessado pelo elevador 2. Eles encontram-se a uma distância possível de ser percorrida a pé, e entre os mirantes há um teleférico que os liga que pode ou não ser utilizado pelo visitante.

O acesso aos elevadores tem os seguintes custos:
\begin{itemize}
\item Subir pelo elevador 1: R\$ 0,15;
\item Subir pelo elevador 2: R\$ 1,80;
\item Descer pelo elevador 1: R\$ 0,10;
\item Descer pelo elevador 2: R\$ 2,30.
\end{itemize}

O custo da passagem do teleférico partindo do topo do mirante 1 para o topo do mirante 2 é de R\$ 2,00, e do topo do mirante 2 para o topo do mirante 1 é de R\$ 2,50.

Qual é o menor custo, em real, para uma pessoa visitar os topos dos dois mirantes e retornar ao solo?
\begin{choices}
\choice 2,25
\choice 3,90
\CorrectChoice 4,35
\choice 4,40
\choice 4,45
\end{choices}

---
tipo: Múltipla Escolha
dificuldade: Média
nivel: médio
disciplina: Matemática
assunto: Porcentagem
gabarito: E
tags: [ENEM, porcentagem, volume, piscina olímpica, capacidade]
fonte: concurso
concurso: ENEM
banca: INEP
ano: 2025
numero: 28
---
\question
O recinto das provas de natação olímpica utiliza a mais avançada tecnologia para proporcionar aos nadadores condições ideais. Para conseguir isso, a piscina de competição tem uma profundidade uniforme de 3 m que ajuda a diminuir a ``reflexão'' da água, além dos já tradicionais 50 m de comprimento e 25 m de largura. Um clube deseja reformar sua piscina de 50 m de comprimento, 20 m de largura e 2 m de profundidade de forma que passe a ter as mesmas dimensões das piscinas olímpicas.

Após a reforma, a capacidade dessa piscina superará a capacidade da piscina original em um valor mais próximo de:
\begin{choices}
\choice 20\%.
\choice 25\%.
\choice 47\%.
\choice 50\%.
\CorrectChoice 88\%.
\end{choices}

---
tipo: Múltipla Escolha
dificuldade: Média
nivel: médio
disciplina: Matemática
assunto: Porcentagem
gabarito: A
tags: [ENEM, porcentagem, desperdício, alimentos, fração]
fonte: concurso
concurso: ENEM
banca: INEP
ano: 2016
cargo: PPL
numero: 29
---
\question
O Brasil é o quarto produtor mundial de alimentos e é também um dos campeões mundiais de desperdício. São produzidas por ano, aproximadamente, 150 milhões de toneladas de alimentos e, desse total, 2/3 são produtos de plantio. Em relação ao que se planta, 64\% são perdidos ao longo da cadeia produtiva (20\% perdidos na colheita, 8\% no transporte e armazenamento, 15\% na indústria de processamento, 1\% no varejo e o restante no processamento culinário e hábitos alimentares).

O desperdício durante o processamento culinário e hábitos alimentares, em milhão de tonelada, é igual a
\begin{choices}
\CorrectChoice 20.
\choice 30.
\choice 56.
\choice 64.
\choice 96.
\end{choices}

---
tipo: Múltipla Escolha
dificuldade: Média
nivel: médio
disciplina: Matemática
assunto: Porcentagem
gabarito: C
tags: [ENEM, porcentagem, lucro, picolé, acréscimo]
fonte: concurso
concurso: ENEM
banca: INEP
ano: 2025
numero: 30
---
\question
Uma pessoa comercializa picolés. No segundo dia de certo evento ela comprou 4 caixas de picolés, pagando R\$ 16,00 a caixa com 20 picolés para revendê-los no evento. No dia anterior, ela havia comprado a mesma quantidade de picolés, pagando a mesma quantia, e obtendo um lucro de R\$ 40,00 (obtido exclusivamente pela diferença entre o valor de venda e o de compra dos picolés) com a venda de todos os picolés que possuía.

Pesquisando o perfil do público que estará presente no evento, a pessoa avalia que será possível obter um lucro 20\% maior do que o obtido com a venda no primeiro dia do evento.

Para atingir seu objetivo, e supondo que todos os picolés disponíveis foram vendidos no segundo dia, o valor de venda de cada picolé, no segundo dia, deve ser
\begin{choices}
\choice R\$ 0,96.
\choice R\$ 1,00.
\CorrectChoice R\$ 1,40.
\choice R\$ 1,50.
\choice R\$ 1,56.
\end{choices}
